\part{Prateek: reading the paper, and checking its images}

\chapter{Prateek's cheat sheet}

\begin{cheat}{Your files, one line each}
\begin{tabularx}{\linewidth}{>{\ttfamily}l X}
pandoc.py & Talks to pandoc: reads a paper into a JSON tree, converts text between Markdown and LaTeX in batches.\\
ingest.py & Turns the tree into a Manuscript: front matter, then each block tagged (heading, figure, table\dots).\\
frontmatter.py & Decides which top lines are author names and which are affiliations; strips LaTeX to plain text.\\
texsource.py & Reads what pandoc loses from LaTeX source: IEEE author blocks, \textbackslash bibitem, .bib, safe macros.\\
manuscript.py & The data model: Manuscript and Block dataclasses, saved as manuscript.json.\\
images.py & Every image: printed dpi, duplicates (dHash), AI evidence in three layers.\\
\end{tabularx}
\end{cheat}

\begin{center}
\begin{tikzpicture}[node distance=5mm and 6mm]
\node[mod] (r) {pandoc.read};
\node[mod, right=of r] (f) {\_flatten};
\node[mod, right=of f] (fm) {\_front\_matter\\\footnotesize frontmatter.py};
\node[mod, below=of fm] (b) {\_body\\\footnotesize tags every block};
\node[mod, left=of b] (md) {\_to\_markdown};
\node[mod, left=of md] (m) {Manuscript};
\node[mod, below=of m] (img) {images.inspect\\\footnotesize dpi \textperiodcentered{} dHash \textperiodcentered{} AI};
\draw[arr] (r) -- (f); \draw[arr] (f) -- (fm); \draw[arr] (fm) -- (b);
\draw[arr] (b) -- (md); \draw[arr] (md) -- (m); \draw[arr] (m) -- (img);
\end{tikzpicture}
\end{center}

\begin{cheat}{Numbers and facts to remember}
\begin{itemize}
\item pandoc node: \cc{\{"t": type, "c": content\}}. Image path is at \cc{c[2][0]}; attributes at \cc{c[0][2]}.
\item Batching: many texts joined by \cc{PRARUPSEPARATOR}, one pandoc run, split again.
\item Only the first \textbf{30} blocks are searched for the title block.
\item Author name: 2 to 4 capitalised words, no digits, no "@", no organisation word.
\item IEEE column \textbf{3.5 in}; minimum \textbf{300 dpi}; dpi = pixels / printed inches (400 px \textrightarrow{} 117 dpi).
\item dHash: 9\(\times\)8 grey pixels, 64 bits; same picture if at most \textbf{4} bits differ.
\item Chart if the 8 most common colours cover \textbf{90\%} of pixels; then the classifier is skipped.
\item Classifier on 26 real figures: \textbf{4} false alarms, all charts; every real AI image scored 1.00.
\item AI codes: IMG-02 declared (error), IMG-03 traces, IMG-04 classifier \(\geq 0.8\).
\end{itemize}
\end{cheat}

\begin{cheat}{The five questions you will surely get}
\textbf{1. How does Prarup read a Word file?} pandoc converts it to a JSON tree; \cc{ingest.open\_paper} flattens it, finds the title block, then tags every block in \cc{\_body}.\par
\textbf{2. How are captions found?} A regular expression (\cc{\_FIG\_CAPTION}) on the paragraph after an image; "Figure 2:" is stripped because IEEE adds "Fig. 2." itself.\par
\textbf{3. How are authors matched to affiliations?} Superscript marks: "Sharp¹" and "¹ Northern Arizona University" share the key "1" (\cc{read\_authors}).\par
\textbf{4. How do you know an image is AI-generated?} We do not claim to know. We report evidence: a signed C2PA record (proof of origin), generator settings left in the file, or optionally a classifier score, never run on charts.\par
\textbf{5. Why batch pandoc calls?} Starting pandoc once per paragraph would take many seconds; one run with separators takes a fraction of one.
\end{cheat}
